Appendix AB — Proton Torsional Fatigue, Spin Stability, and Mass Potential

Framework: Scarlet–VanAcker Cosmological Framework (v2.0)
Objective: Model the mechanical relaxation of the nucleon knot core as a function of cosmic expansion, determine spin stability, lattice memory, and the contribution of torsional energy to proton mass.

⸻

1. The Fatigue Function: Internal Torsional Pressure

As the 11D filament lattice expands, the mechanical coupling between the VanAcker Bedrock and 4D spacetime undergoes elastic hysteresis. The internal torsional pressure of the nucleon knot is modeled as:
	•	T_internal(z) = T_0 × [1 – ε_f × f(z)]

Where:
	•	T_0 ≈ 1.05 × 10³⁵ Pa is the baseline torsional pressure at z = 0,
	•	ε_f is the Lattice Fatigue Coefficient, representing the mechanical “memory” of spacetime,
	•	f(z) = (1 + z)^(3/2) represents strain scaling with cosmic expansion.

As redshift increases toward z ≈ 1.5, the core torsional gradient softens, but the nucleon remains a topologically stable knot.

⸻

2. Spin Conservation and Topological Invariance

In Scarlet 2.0, the proton’s spin S = 1/2 is a topological winding number of the filament lattice. Fatigue redistributes internal torque, but total angular momentum is conserved:
	•	S = ∫_0^∞ T(r, z) dr = 1/2 ħ

As the torsional radius r_τ expands with fatigue, the knot softens mechanically without changing total spin. Spin is preserved even as the internal torque distribution shifts.

⸻

3. Hyperfine Splitting and Observational Consequences

The redistribution of internal torque affects the proton magnetic moment μ_p, leading to minor shifts in hyperfine transition frequencies:
	•	Relative frequency shift: δ(Δν)/Δν ≈ γ × ε_f × f(z)

Where γ is the torsional compliance constant. At z ≈ 1.5, the relative shift is on the order of 10⁻⁸, potentially observable with high-precision spectroscopy.

Observational probes include:
	•	Quasar absorption spectroscopy (probing μ variations),
	•	21cm neutral hydrogen hyperfine transitions at high redshift,
	•	Atomic clock comparisons sensitive to nuclear structure variations.

⸻

4. Mechanical Interpretation and Lattice Memory

Proton fatigue reflects a gradual redistribution of torsional stress. Key points:
	•	The torsional lattice behaves like an elastic medium with extremely high memory.
	•	If cosmic expansion were to halt, partial re-stiffening could occur, but full relaxation is a thermodynamically irreversible process tied to entropy increase in the lattice.
	•	ε_f is extremely small (~10⁻¹² yr⁻¹ for terrestrial conditions), meaning relaxation is negligible at z = 0 but accumulates over gigayear timescales.

⸻

5. Cosmic Evolution of the Proton Knot
	•	Early Universe (z > 10): Lattice stiffness G_τ is extremely high. Proton knots are “super-stiff,” torsional energy storage is maximal, and fatigue is effectively zero.
	•	Transition Era (z ≈ 1.5): Fatigue becomes measurable. Internal torque redistributes, r_τ slightly expands, and hyperfine splitting shifts minimally.
	•	Modern Era (z = 0): Proton knots are slightly softer, r_τ ≈ 0.073 fm, and torsional pressure reaches equilibrium (~1.05 × 10³⁵ Pa).

The torsional “Hard Core” of the nucleon was smaller in the early universe (r_τ ≈ 0.062 fm at z ≈ 10), implying higher internal pressures and a more compact nucleon core. This could affect neutron star equations of state for high-redshift remnants (“Scarlet Correction”).

⸻

6. Torsional Mass Potential

The nucleon knot stores mechanical energy in the filament lattice. This torsional energy contributes a small fraction of proton mass:
	•	E_torsion ≈ 0.5 × P_peak × V_eff ≈ 1.25 × 10⁻¹⁰ J
	•	Corresponding mass contribution: Δm_torsion ≈ 0.78 MeV/c² (~0.1% of proton mass)

Fatigue slightly modifies Δm_torsion over cosmic time, but spin and topological stability remain intact. Early-universe nucleons had a larger torsional energy contribution due to higher lattice stiffness, explaining a “super-stiff” primordial plasma.

⸻

7. Quantitative Constraints and Predictive Mapping
	•	VanAcker Constant Constraint: ε_f < 10⁻¹² yr⁻¹ for terrestrial measurements.
	•	Effective torsional modulus evolves as: G_τ(z) = G_τ(0) × [1 + η × ln(1+z)], with η ≈ 10⁻⁴ capturing early-universe hardening.
	•	Fatigue effects are cumulative but reversible only partially under high compression.

⸻

8. Technical Implications
	•	Proton fatigue links microphysical torsion to large-scale cosmic evolution.
	•	The torsional radius and internal energy distribution affect nuclear physics and neutron star EoS modeling at high redshift.
	•	The lattice memory and torsional mass potential provide a predictive framework for subtle variations in fundamental constants across cosmological time.

⸻

9. Summary
	•	Proton spin remains exactly 1/2 throughout cosmic expansion.
	•	Torsional fatigue redistributes internal stress without breaking topological continuity.
	•	Early-universe nucleons were more compact, “super-stiff,” and stored more torsional energy.
	•	Modern protons exhibit a slightly softer knot structure with minor mass and magnetic moment shifts.
	•	Observables: hyperfine transitions, quasar absorption lines, atomic clock comparisons, and neutron star EoS constraints.
